% De Re Publica (de-re-publica) --- English translation
% Translated by Claude (Anthropic) from the Latin (Perseus).
% Style guide: STYLE.md.

\ciceroBook{1}

\ciceroSection{1} \ldots had not, by their onset, freed [the state]; nor had Gaius Duilius, Aulus Atilius, Lucius Metellus delivered Rome from the terror of Carthage; nor had the two Scipios put out, with their own blood, the rising flame of the Second Punic War; nor, when it broke out again with greater forces, would either Quintus Maximus have worn it down, or Marcus Marcellus crushed it, or Publius Africanus, having torn it from the very gates of this city, driven it back within the enemy's walls. To Marcus Cato, an obscure new man --- whom all of us who pursue the same things are drawn after, as if to a model, toward industry and virtue --- it was certainly open to enjoy himself in leisure at Tusculum, a healthful spot and close at hand. But the man, mad (as those people think), with no necessity compelling him, preferred to be tossed in these waves and storms into extreme old age rather than to live, in such tranquillity and leisure, most agreeably. I pass over the countless men, each of whom was a saving of this commonwealth: because they are not far off from the memory of this age, I leave their names alone, lest someone complain that he himself, or one of his own, has been overlooked. This one thing I lay down: that so great is the necessity of virtue given by nature to the human race, and so great the love of defending the common safety, that this force has overcome all the enticements of pleasure and ease.

\ciceroSection{2} Nor indeed is it enough to possess virtue as if it were some art, unless you use it; for although an art, even when not in use, can still be held by knowledge alone, virtue lies wholly in its exercise; and the greatest exercise of virtue is the steering of the state and the carrying out --- in fact, not in talk --- of those very things which those people declaim about in corners. For nothing is said by the philosophers, at least when it is rightly and honourably said, that has not been begotten and confirmed by those by whom legal codes were drawn up for cities. Whence comes piety, or from whom religion? Whence comes law --- whether of nations, or this very thing called civil law? Whence justice, faith, equity? Whence shame, restraint, the flight from baseness, the appetite for praise and honour? Whence in toils and dangers comes courage? Surely from those who, having shaped these things by training, have confirmed some by custom and sanctioned others by law.

\ciceroSection{3} Indeed they say of Xenocrates, the foremost philosopher of his rank, that when he was asked what his pupils gained, he answered: ``That they should do of their own will what they would be compelled to do by the laws.'' That citizen, then, who by command and by penalty of the laws compels everyone to do what philosophers can scarcely persuade a few to do by their speech, is to be set above the very teachers who debate the matter. For what speech of theirs is so refined that it is to be preferred to a well-ordered state with its public law and its customs? Just as I think ``great and imperious cities,'' as Ennius calls them, are to be set before villages and forts, so I judge those who preside over these cities by counsel and authority to be set far above, in wisdom itself, those who have no part in any public business. And since we are most of all swept on toward enlarging the resources of the human race, and study to render the life of men safer and richer by our counsels and our labours, and since we are spurred to this very pleasure by the goads of nature herself, let us hold to that course which has always been the course of the best men, and let us not listen to those signals which sound the recall, calling back even those who have already gone forward.

\ciceroSection{4} Against these reasonings, so sure and so brilliant, those who argue on the other side set first the labours which must be borne in defending the commonwealth --- a slight enough hindrance, surely, to the watchful and energetic man, and to be despised not only in such great matters but even in middling pursuits or duties or even in business affairs. Then they add the dangers of life, and they hold up before brave men a base fear of death --- to whom this is wont to seem more wretched, that they should be used up by nature and old age, than that the time should be granted them to give back, on behalf of their country above all, that life which had to be given back to nature in any case. And it is on this ground that they think themselves rich and eloquent, when they gather up the calamities of the most distinguished men and the wrongs done them by their ungrateful fellow citizens.

\ciceroSection{5} For from this come those examples both among the Greeks: that Miltiades, the conqueror and tamer of the Persians, when the wounds were not yet healed which he had taken in his glorious victory by his front turned to the enemy, poured out, in the chains of his own citizens, the life saved from the enemy's weapons; and that Themistocles, driven from the country he had freed and frightened off, took refuge not in the Greek harbours he himself had preserved, but in the gulfs of barbarism, which he had crushed. Nor are examples wanting of Athenian fickleness and cruelty toward their best citizens; and those things, born and rife with them, are said to have spilled over even into our own most weighty state ---

\ciceroSection{6} for they bring up either Camillus's exile, or the affront to Ahala, or the unpopularity of Nasica, or the expulsion of Laenas, or Opimius's condemnation, or Metellus's flight, or the most bitter ruin of Gaius Marius and the slaughter of leading men, or the destruction of those many others which followed not long after. And they no longer keep off my own name; and, I suppose, because they think themselves preserved by my counsel and at my danger to that life and leisure of theirs, they complain of me with even more weight, and more affection. But I cannot easily say why, when these very men cross the seas for the sake of learning and seeing \ldots

\ciceroSection{7} \ldots [had I, swearing] that the commonwealth was safe, when leaving the consulship in the assembly, with the Roman people swearing the same, I should easily compensate the trouble and vexation of all the wrongs done me. Yet our own misfortunes brought more honour than toil, and not so much vexation as glory, and we drew greater joy from the longing of the good than pain from the joy of the wicked. But if it had turned out otherwise, as I have said, how could I have complained? When nothing had befallen me unexpectedly nor more grievously than I had looked for, in return for such great deeds of mine. For I had been such a man that, although it was open to me either to take greater fruits from leisure than others, by reason of the varied sweetness of the studies in which I had lived from boyhood, or, if anything more bitter befell us all, to undergo not a special, but only an equal share with the rest in fortune, I did not hesitate to throw myself headlong against the heaviest storms and almost the very thunderbolts themselves, for the sake of saving my fellow citizens, and by my own personal dangers to procure the common quiet for the rest.

\ciceroSection{8} For our country did not beget us or rear us under such a law that she should expect, as it were, no nourishment from us, and serve only our conveniences, supplying a safe refuge for our leisure and a peaceful place for our quiet --- but rather that she should pledge to herself, for her own use, the most and the largest part of our spirit, our talent, our judgement, and grant back to us, for our private use, only as much as could be left over for herself.

\ciceroSection{9} Then those refuges they take to themselves for an excuse, that they may the more easily enjoy their leisure --- those, too, are by no means to be listened to, when they say that men who come to public affairs are for the most part fellows worthy of nothing good, with whom to be compared is sordid, and to fight, especially when the multitude is roused, is wretched and dangerous. For which reason, they say, it is neither the wise man's part to take up the reins, when he cannot restrain the mad and untamed onsets of the rabble, nor the part of a free man, when he is contending with foul and monstrous opponents, to bear the lashes of insults or to expect injuries which the wise man cannot endure --- as if good men, brave men, men endowed with great spirit had any juster cause for entering on public business than this: not to obey the worthless, nor to suffer the commonwealth to be torn apart by them, when they themselves, even if they wished to bring help, could not.

\ciceroSection{10} And that exception of theirs --- to whom can it ever be approved? --- when they say the wise man will undertake no part in the commonwealth except when time and necessity have compelled him: as if a greater necessity could befall any man than has befallen me. In that necessity, what could I have done if I had not been then consul? But how could I have been consul, if I had not held that course of life from boyhood by which, born of equestrian rank, I might come to the most distinguished honour? There is no power, then, to bring help to the commonwealth on the spur of the moment, or when you wish, however hard pressed it is by dangers, unless you are in the place where you may do so.

\ciceroSection{11} And what most often seems strange to me in the speech of these learned men is this: that those who deny they can steer in a calm sea --- because they have neither learned the art nor ever cared to know it --- the very same men profess that they will set hand to the helm when the greatest waves have been raised. For these people are wont to say openly, and even to glory in the saying, that they have neither ever learned nor teach anything about the principles of either founding or preserving states; and they think the knowledge of these things is to be conceded not to learned and wise men, but to those skilled in that other line. How then is it fitting to promise their service to the commonwealth only when they are forced by necessity --- since, what is much easier, when no necessity presses, they cannot govern the state? For my part, even if it were true that the wise man does not commonly come down of his own will to public business, but if compelled by the times then does not refuse the duty, I should still judge that this knowledge of public affairs is by no means to be neglected by the wise man, since he ought to have everything ready for use, when he might not know whether at some time he must use it.

\ciceroSection{12} I have spoken of these matters at greater length for this reason: that in these books a discussion has been undertaken and taken on by me concerning the commonwealth; and lest it be conducted in vain, I had first of all to remove this hesitation about entering on public business. Yet if there are any who are moved by the authority of philosophers, let them give heed for a little while and listen to those whose authority and glory is highest among the most learned men --- whom I judge, even if some of them did not themselves manage public business, still, since they have inquired and written much about the commonwealth, to have discharged some duty to the commonwealth. As to those seven whom the Greeks named the Wise Men, I see that almost all of them moved in the very midst of public business. For there is no thing in which human virtue draws nearer to the divine power of the gods than in either founding new cities or preserving those already founded.

\ciceroSection{13} On these matters, since it has been our lot both to have achieved something worthy of remembrance in the conduct of the commonwealth and to be authorities, by some skill in the unfolding of the principles of public affairs, not only by use but also by the study of learning and teaching --- whereas earlier men were either polished in their disputations, but no deeds of theirs are to be found, or excellent in action but unpracticed in argument --- still no new doctrine of our own, invented by us, is to be set up here: rather, we must call back to memory the discussion of the most distinguished and wisest men of our state, of one age, which was reported to me, and to you when you were a young man, by Publius Rutilius Rufus when we were together at Smyrna for several days; in which scarcely anything was passed over which had any great bearing on the principles of the whole subject.

\ciceroSection{14} For when this Publius Africanus, son of Paulus, in the consulship of Tuditanus and Aquilius, had decided to be in his gardens for the Latin holidays, and his closest friends had said they would visit him often during those days, on the morning of the holidays itself the first to come to him was the son of his sister, Quintus Tubero. When Scipio had hailed him kindly and was glad to see him, he said: ``Why so early, Tubero? These holidays were giving you a fine opportunity for unfolding your books.'' Then Tubero: ``For me indeed every time is free for my books; for they are never busy. But to find you at leisure is no small thing, especially with this present commotion of the commonwealth.'' Then Scipio: ``You have caught me, but, by Hercules, more at leisure in business than in mind.'' And Tubero: ``Yet you ought to relax your mind too. We are many, as we determined, ready, if it can be done at your convenience, to enjoy this leisure with you.'' ``Gladly will I --- so that we may at last be put in mind of something out of the studies of learning.''

\ciceroSection{15} Then he said: ``Do you wish, then, since you in some way invite me and give me hope of yourself, that we look first into this, Africanus, before the others come --- what is the meaning of that other sun reported in the senate? For neither few nor light are those who say they have seen two suns, so that the matter is less one of crediting than of looking for the cause.'' Here Scipio: ``How I should wish we had our Panaetius with us! He, with the rest, is wont to investigate these celestial things most eagerly. But for my part, Tubero --- to speak openly to you what I think --- I do not greatly agree, in this whole subject, with that intimate of ours, who, when we can scarcely guess by conjecture what these things are, so confidently affirms them as if he saw them with his eyes or fully handled them with his hand. For this reason I am wont to judge Socrates the wiser, who set aside all care of this kind, and said that the things asked about nature were either greater than the human mind could attain, or had no bearing at all on the life of men.''


